% f4_space.tex -- Figure 4: the profile space, drawn to scale for CIFAR-10 (natural width about 17.3 cm).
%   Whitened coordinates: one unit = one ID standard deviation measured with P (P ~ S^-1), so the ordinary distance in the picture is the
%   Mahalanobis distance.  Horizontal axis: the direction of mu - mu_o; vertical: one of the other p - 1 = 22 directions.
%   mu = (0, 0), mu_o = (-Delta, 0), mu* = (c Delta, 0) with Delta = 1.768 and c = 8 (olx_info.json).  The clouds are schematic
%   (gen/mk_data.py).  Red dashed circles: equal score (circles about mu*).  Colours: navy = ID, grey = re-oriented, red = mu* and OOD.
\begin{tikzpicture}[>={Stealth[length=1.8mm, width=1.4mm]}, line cap=round, line join=round, font=\footnotesize,
  lab/.style={inner sep=1pt, align=center, font=\footnotesize}, sm/.style={inner sep=1pt, align=center, font=\scriptsize}]
% written by gen/mk_data.py: schematic clouds in whitened coordinates (unit = one ID standard deviation), Delta = 1.768
\def\iddots{1.141/-0.915, 0.793/-0.396, -0.460/0.787, 0.258/-0.211, -0.391/-0.483, 1.610/-0.506, 0.348/-0.784, 1.777/0.323, -2.663/0.210, -1.614/1.601, 0.070/-0.577, -0.432/-0.972, 0.960/-1.231, -0.113/-1.618, -0.096/-0.026, -0.013/-0.091, 1.323/0.224, 0.288/0.800, 1.522/-0.368, -0.714/0.926, -1.364/0.934, 1.369/0.839, 1.282/-2.040, -0.463/-0.773, 0.823/-0.255, 0.763/-0.393, 1.487/-1.352, 1.386/0.398, -0.270/0.410, 0.430/0.463, 0.877/-0.427, -0.372/-0.239, 1.049/-0.658, 0.223/-1.839, -0.944/-0.816, 0.348/0.524, -1.523/0.989, 2.220/-0.857, -1.523/1.908, -0.124/-0.016, -0.895/0.780, -1.154/-0.740, -0.883/-0.517}
\def\ordots{-1.181/-0.324, -2.756/0.911, 0.444/0.636, -3.820/-0.668, -1.408/0.078, -1.856/1.755, -0.855/0.710, -3.603/0.487, -2.943/1.209, -1.415/-1.033, -1.930/0.523, -0.329/-0.267, -2.925/-0.010, -2.026/0.355, -2.910/-0.731, -2.447/-0.342, -1.286/1.351, -3.754/-0.026, -1.556/0.315, -3.228/-0.173, -2.023/0.231, -2.999/-0.910, 0.621/0.860, -1.495/0.729, -2.527/0.524, -0.797/0.393, -0.870/0.235, -1.771/0.334, -2.710/-1.838, -3.078/-1.015, -1.412/0.120, -2.912/0.736, -2.381/2.031, -1.829/0.285, -2.386/0.150, -2.385/1.099, -1.830/0.783, -2.075/-0.558, -1.726/1.890, -2.952/-0.393, -1.908/-0.029, -2.152/-1.155, 0.660/-0.254, -2.828/0.498}
\def\mDelta{1.768}

\begin{scope}[shift={(3.75,0)}, x=0.78cm, y=0.78cm]
  \pgfmathsetmacro{\cD}{8*\mDelta}
  \draw[black!35, line width=0.4pt] (-4.6,-2.55) rectangle (15.3,2.55);
  \begin{scope}
    \clip (-4.6,-2.55) rectangle (15.3,2.55);
    \draw[black!25, line width=0.4pt] (-4.6,0) -- (15.3,0);
    \foreach \r in {3,6,...,18} \draw[redc!75, densely dashed, line width=0.5pt] (\cD,0) circle (\r);
    \foreach \a/\b in \ordots \fill[black!40] (\a,\b) circle (1.25pt);
    \foreach \a/\b in \iddots \fill[navyc] (\a,\b) circle (1.25pt);
  \end{scope}
  \draw[navyc, line width=0.5pt] (0,0) circle (1);
  \node[sm, text=navyc, anchor=north, fill=white, inner sep=0.5pt] at (0,-1.06) {1標準偏差};
  % means and the two arrows
  \fill[black!60] (-\mDelta,0) circle (2.2pt);
  \fill[navyc] (0,0) circle (2.2pt);
  \draw[->, line width=1.0pt] (-\mDelta+0.12,0) -- (-0.12,0);
  \node[lab, anchor=south, fill=white, inner sep=0.5pt] at (-0.88,0.16) {$\vec{\mu}-\vec{\mu}_{o}$};
  \node[lab, anchor=north, fill=white, inner sep=0.5pt] at (-\mDelta,-0.15) {$\vec{\mu}_{o}$};
  \node[lab, anchor=north, text=navyc, fill=white, inner sep=0.5pt] at (0.05,-0.15) {$\vec{\mu}$};
  \draw[->, redc, densely dashed, line width=1.0pt] (0.15,0) -- (\cD-0.2,0);
  \node[lab, text=redc, anchor=south, fill=white, inner sep=1pt] at (6.95,0.12) {$c(\vec{\mu}-\vec{\mu}_{o})$, \ $c=8$};
  \node[redc, font=\Large, inner sep=0] at (\cD,0) {$\star$};
  \node[lab, text=redc, anchor=south, fill=white, inner sep=1pt] at (\cD,0.3) {$\vec{\mu}^{*}$(基準点)};
  % the two test images
  \draw[navyc, line width=0.9pt] (0.5,0.8) circle (3.2pt);
  \draw[navyc, line width=0.4pt] (0.62,0.95) -- (1.45,1.75);
  \node[sm, text=navyc, anchor=west, fill=white, inner sep=1pt] at (1.45,1.82) {分布内のテスト画像};
  \fill[redc] (-3.3,-0.8) circle (3.0pt);
  \draw[redc, line width=0.4pt] (-3.3,-0.95) -- (-3.2,-1.35);
  \node[sm, text=redc, anchor=north, fill=white, inner sep=1pt, align=center] at (-3.1,-1.35) {分布外のテスト画像\\(振幅が下がった)};
  % names of the clouds
  \node[sm, text=black!65, anchor=north west, fill=white, inner sep=1pt] at (-4.5,2.45) {向きを変えた画像};
  \node[sm, text=navyc, anchor=west, fill=white, inner sep=1pt] at (1.9,-1.55) {分布内の画像};
  \node[sm, text=redc, anchor=north, fill=white, inner sep=1pt, align=center] at (7.0,-1.4) {スコアが等しい線\\($\vec{\mu}^{*}$ を中心とする円)};
  % axes
  \node[sm, anchor=west, align=left] at (15.35,0) {$\vec{\mu}-\vec{\mu}_{o}$\\の向き};
  \node[sm, anchor=south west] at (-4.6,2.58) {縦：ほかの向き(実際は $p-1=22$ 個). 縮尺：分布内の標準偏差をそろえた座標(ここでの距離がマハラノビス距離)};
  \draw[<->, line width=0.6pt] (-2.0,-3.0) -- (10.5,-3.0);
  \node[sm, anchor=east] at (-2.1,-3.0) {スコアが低い};
  \node[sm, anchor=west] at (10.6,-3.0) {スコアが高い};
\end{scope}
\end{tikzpicture}
